\section{A realized three-dimensional family with exact work counts}
\label{fib:section}
\label{fib:family}

The following family separates the geometric advantage of minimum cells
from questions about which sectors an unknot recognizer must visit.  Every
member is an actual triangulated solid torus.  Its supplied sector has
three-dimensional quadrilateral matching space, seven nonlink standard
rays, and exponentially large primitive coordinates.  All statements
quantified over the family are proved below; the finite computational
checks are reported separately at the end.

\subsection{Face pairings and the underlying manifold}
\label{fib:source}

Number tetrahedron vertices \(0,1,2,3\), and number a face by its opposite
vertex.  A permutation written \((p_0,p_1,p_2,p_3)\) sends vertex \(j\) to
\(p_j\).  A normal row is ordered
\[
 (t_0,t_1,t_2,t_3;q_0,q_1,q_2),
\]
where quadrilateral types \(0,1,2\) avoid, respectively, the opposite
edge pairs \((01,23)\), \((02,13)\), and \((03,12)\).
These are the conventions of the native fixture.

For \(n\geq1\), let \(T_n\) contain tetrahedra \(0,\ldots,n-1\), with
face pairings
\begin{align}
 (i,0)&\xrightarrow{(2,1,3,0)}(i+1,2),&
 (i,1)&\xrightarrow{(0,3,1,2)}(i+1,3)
       &&(0\leq i<n-1),\label{fib:layer-pairings}\\
 (n-1,0)&\xrightarrow{(1,2,3,0)}(n-1,1).
 \label{fib:core-pairing}
\end{align}
The inverse pairings are understood.  Its unpaired faces are \((0,2)\)
and \((0,3)\).  Form \(M_n\) by adding tetrahedra \(n,n+1\) and pairing
\begin{equation}
 (0,2)\xrightarrow{\mathrm{id}}(n,2),\qquad
 (0,3)\xrightarrow{\mathrm{id}}(n+1,3).
 \label{fib:cap-pairings}
\end{equation}
Thus \(t=n+2\).  Allow type \(2\) in each original tetrahedron and type
\(1\) in each cap, so \(k=t\).

The face closures in this description need not embed: all vertices of
\(T_n\) are identified.  The following observation handles that issue
without treating a quotient face as an embedded closed disc.

\begin{lemma}[Boundary caps with identified outer boundaries]
\label{fib:collar}
Let \(S=\partial M\) be the boundary of a compact PL manifold.  Suppose
an abstract PL disc \(D\) is attached by a characteristic map
\(f:D\to S\) that embeds its interior, with
\(f(\partial D)\cap f(\operatorname{int}D)=\varnothing\).
Identifications between points of \(\partial D\) are permitted.
Attach a PL ball along a disc in its boundary identified with \(D\).
The resulting space is PL homeomorphic to \(M\).  Its boundary replaces
\(D\) by the complementary abstract boundary disc, with the same outer
boundary identifications.
\end{lemma}
\begin{proof}
Choose a PL function \(\rho:D\to[0,1]\), positive in the interior and
zero on the boundary, after subdividing \(D\) if necessary.  The region
\[
 B_\rho=\{(p,s):p\in D,\ 0\leq s\leq\rho(p)\}
\]
is a PL ball whose lower and upper discs meet along \(\partial D\).
A boundary homeomorphism from the given ball to this model, fixed on
the attaching disc, extends across the balls by coning a common boundary
subdivision.  Since every identification under \(f\) occurs where
\(\rho=0\), the height function descends to a continuous PL function
\(\bar\rho\) on \(S\), extended by zero off \(f(D)\).
The pushout is therefore the original manifold with the region
\(0\leq s\leq\bar\rho(q)\) added to an exterior boundary collar.

Take an interior collar \(S\times[-1,0]\).  The enlarged collar has
vertical intervals \([-1,\bar\rho(q)]\).  Triangulate \(S\) so that
\(\bar\rho\) is affine on each simplex.  Corresponding triangulations
of the prisms between the lower and upper graphs give a PL homeomorphism
to \(S\times[-1,0]\), fixing the lower boundary.  Extend it by the
identity outside the collar.  This also proves the assertion at identified
vertices and edges: it is an ambient collar construction, not merely a
comparison of the two open faces.
\end{proof}

\begin{proposition}[Realization as solid tori]
\label{fib:solid-torus}
The space \(T_n\) is an orientable solid torus with one vertex.  In its
front boundary tetrahedron its three distinct boundary-edge classes are
\[
 a_n=[01],\qquad b_n=[02]=[13],\qquad c_n=[03]=[12].
\]
The space \(M_n\) is an orientable solid torus with three boundary
vertices.
\end{proposition}
\begin{proof}
For \(T_1\), the edge classes are
\(\{01\}\), \(\{02,13\}\), and \(\{03,12,23\}\), all boundary
edges.  Its four vertex-link triangles are joined in the chain
\(1,2,3,0\), giving one disc link; its edge links are intervals, without
reversed self-identifications.  The face pairing is orientation compatible.
Thus this is a compact orientable manifold.  Removing collars of its
boundary vertices, edges and faces leaves a ball around the tetrahedron
centre and one handle through the paired faces.  There are no interior
edge or vertex dual blocks.  Consequently the manifold is a regular
neighbourhood of a one-loop graph, hence an orientable solid torus.

To pass from \(T_m\) to \(T_{m+1}\), put a new front tetrahedron before
the old ones, using~\eqref{fib:layer-pairings}.  Its attaching disc is
the union of faces \(0,1\), with internal diagonal \(23\).  Both
occurrences of this diagonal map to \(c_m\); the four outer edges map
to \(a_m,b_m,a_m,b_m\).  These three edge classes are distinct.
The interior of the abstract attaching disc therefore embeds in the old
boundary: its only boundary identifications concern the outer edges and
vertices.  Lemma~\ref{fib:collar} applies.  The new boundary classes are
\(a_{m+1}=[01]\), \(b_{m+1}=a_m\), and \(c_{m+1}=b_m\);
the old \(c_m\) becomes interior.  The new vertices join the unique
old vertex.  This proves the induction and the stated distinctness.

Each attachment in~\eqref{fib:cap-pairings} satisfies the one-face case
of the same lemma.  The remaining original face is unaffected by the
first attachment.  Each cap contributes its opposite vertex as a new
boundary vertex.  In link language, the three old corner intervals are
distinct intervals on the old vertex-link boundary; the cap adds triangular
flaps there, while its new vertex has a disc link.  These observations
also explain directly why the quotient face closures cause no singularity.
\end{proof}

\subsection{The complete matching cone and its canonical section}
\label{fib:matching}

Set \(F_0=0\), \(F_1=1\), and \(F_{j+2}=F_{j+1}+F_j\), and abbreviate
\(A=F_{n+1}\), \(B=F_{n+2}\).

\begin{theorem}[Full coordinates, including vertex links]
\label{fib:coordinates}
Every real standard matching solution in the prescribed support has the
following unique form, with parameters \(x,y,z,b,u,v\):
\begin{align}
 \text{row }i:\quad&
 (F_{n-i+1}x+b,F_{n-i+1}x+b,b,b;0,0,F_{n-i}x),
 &&0\leq i<n,\label{fib:base-row}\\
 \text{row }n:\quad&
 (Ax+b-y,Bx+b,u,b;0,y,0),\label{fib:first-cap}\\
 \text{row }n+1:\quad&
 (Bx+b,Ax+b-z,b,v;0,z,0).\label{fib:second-cap}
\end{align}
It is nonnegative precisely when
\begin{equation}
 x,b,u,v\geq0,\qquad 0\leq y,z\leq Ax+b.
 \label{fib:full-box}
\end{equation}
Its quadrilateral matching space has dimension three; its quadrilateral
cone is parametrized by \((x,y,z)\in\R_{\geq0}^3\).
For such a triple put
\begin{equation}
 h=\max(0,y-Ax,z-Ax).
 \label{fib:hinge}
\end{equation}
The canonical normal vector has \(b=h\) and \(u=v=0\).
Every nonnegative solution uniquely adds \((b-h)\) copies of the old
vertex link and \(u,v\) copies of the two new vertex links to that
canonical vector.
\end{theorem}
\begin{proof}
Substitute type \(2\) into the base face equations, working backwards
from~\eqref{fib:core-pairing}.  The triangle entries are
\((X_i,X_i,Y_i,Y_i)\), with quadrilateral entry \(W_i\), and the
equations reduce to
\begin{align*}
 X_i&=X_{i+1}+W_{i+1},&Y_i&=Y_{i+1},&
 Y_i+W_i&=X_{i+1}\quad(0\leq i<n-1),\\
 X_{n-1}&=Y_{n-1}+W_{n-1}.
\end{align*}
Taking \(x=W_{n-1}\) and \(b=Y_{n-1}\) gives
\(W_i=F_{n-i}x\), \(X_i=F_{n-i+1}x+b\), and \(Y_i=b\).
Conversely these formulas satisfy every base equation.
Writing the first cap triangles as \(r_j\), its attaching equations are
\(r_0+y=Ax+b\), \(r_1=Bx+b\), and \(r_3=b\); \(r_2=u\)
is free.  For the second cap they are
\(r_0=Bx+b\), \(r_1+z=Ax+b\), and \(r_2=b\); \(r_3=v\)
is free.  This proves the complete linear parametrization and
\eqref{fib:full-box}.

Every nonnegative \((x,y,z)\) has a lift by choosing \(b\geq h\).
The last base quadrilateral and the two cap quadrilaterals recover these
three parameters, so the parametrization is injective and its matching
dimension is three.  At the old global vertex the minimum triangle entry
is
\[
 \min(b,Ax+b-y,Ax+b-z)=b-h;
\]
all other old-vertex entries are at least \(b\).  The two new minima
are \(u,v\).  Subtracting these three minima proves the canonical and
link assertions, in the standard sense of canonical
quadrilateral-to-standard lifting~\cite{Burton2009}.
\end{proof}

\begin{proposition}[Seven nonlink rays and three active cells]
\label{fib:seven}
The complete standard sector has ten primitive rays: its three vertex
links and the canonical lifts of
\begin{equation}
 \begin{split}
 &(1,0,0),\ (1,A,0),\ (1,0,A),\ (1,A,A),\\
 &(0,1,0),\ (0,0,1),\ (0,1,1).
 \end{split}
 \label{fib:seven-parameters}
\end{equation}
The normalized quadrilateral section has three minimum cells, three
internal edges and seven vertices.  The one-group output bound is attained:
\(3+2(3-1)=7\).
\end{proposition}
\begin{proof}
First discard the independent coordinates \(u,v\).  Normalize the
remaining cone by \(Ax+b=1\).  Equation~\eqref{fib:full-box} becomes
the cube
\[
 (Ax,y,z)\in[0,1]^3,\qquad b=1-Ax.
\]
Its eight vertices give four rays with \(x=1,b=0\) and
\(y,z\in\{0,A\}\), and four with \(x=0,b=1\) and
\(y,z\in\{0,1\}\).  The latter corner \(y=z=0\) is the old
vertex link; the other three are the last three rays in
\eqref{fib:seven-parameters}.  Restoring \(u,v\) adds the two other
link rays.  The displayed nonlink vectors are primitive because either
the last base quadrilateral or a cap quadrilateral is one.

In the base-offset gauge the only active minimum functions are
\(0,Ax-y,Ax-z\).  Their three regions are
\[
 y,z\leq Ax;\qquad y\geq Ax,\ y\geq z;\qquad
 z\geq Ax,\ z\geq y.
\]
The actual normalization is \((B-1)x+y+z=1\), since
\(\sum_{j=1}^nF_j=B-1\).  On this triangle the three internal edges
meet at the direction \((1,A,A)\) and end at
\((1,A,0)\), \((1,0,A)\), and \((0,1,1)\).
Together with the original triangle vertices these are exactly seven
vertices, proving the geometric assertions.
\end{proof}

\subsection{A complete disc census for every integral parameter}
\label{fib:topology}

\begin{lemma}[The base meridian]
\label{fib:meridian}
The base vector \eqref{fib:base-row} with \(x=1,b=0\) is one essential
meridian disc \(D_n\).
\end{lemma}
\begin{proof}
Any normal component has nonnegative integral base parameters.  A
decomposition of this vector has total \(x=1\), \(b=0\), so exactly
one component can be nonzero.  This proves connectedness without expanding
the Fibonacci coordinates.

Here is an explicit Euler-characteristic check.  Write
\(C=F_{n+3}=A+B\).  The three boundary-edge intersection numbers, in
the order \(a_n,b_n,c_n\), are \(C,B,A\).
The interior edges inherited from \(T_m\), \(1\leq m<n\), have
weights \(F_{m+1}\).  Thus the normal cell decomposition has
\[
 V=2C+B-2,\qquad E=4C+2B-8,\qquad F=2C+B-5.
\]
Indeed there are \(2(C-2)\) triangles and \(B-1\) quadrilaterals;
their sides, with boundary sides counted once, give the displayed edge
count.  The vertex count follows by summing the edge weights and using
\(\sum_{m=1}^{n-1}F_{m+1}=B-2\).  Hence \(V-E+F=1\).
The connected surface has nonempty boundary and Euler characteristic one,
so it is a disc.

An inessential circle on the boundary torus has even mod-two intersection
with every closed edge loop.  Consecutive Fibonacci numbers \(A,B\)
are coprime, so they cannot both be even.  The listed boundary intersections
therefore show that \(\partial D_n\) is essential.  In a solid torus
an essential properly embedded disc is a meridian disc.
\end{proof}

\begin{theorem}[Integral component classification]
\label{fib:disc-census}
For nonnegative integral \((x,y,z)\), the canonical vector consists of
\(x\) essential meridian discs and \(h\) inessential discs, where
\(h\) is given by~\eqref{fib:hinge}.  Its component count and Euler
characteristic are both \(x+h\).  More generally, the full integral
solution~\eqref{fib:base-row}--\eqref{fib:second-cap} consists of
\(x\) meridian discs and \(b+u+v\) inessential discs.
\end{theorem}
\begin{proof}
First take \(u=v=0\).  The base restriction is the disjoint union of
\(x\) parallel copies of \(D_n\) and \(b\) small vertex-link discs.
These can be chosen disjoint: the meridian avoids the vertex, and its
two-sided regular neighbourhood supplies the parallel copies.  Their
normal coordinates are exactly the base rows.  The usual uniqueness of
an embedded normal surface with specified coordinates identifies this
union with the base restriction: stack the prescribed pieces in each
tetrahedron and match the ordered normal arcs across its faces.

Every cap piece present when \(u=v=0\) meets its attaching face in
exactly one arc.  This holds for each quadrilateral and for each triangle
whose vertex belongs to that face.  The sole triangle type missing the
face is its opposite-vertex triangle, and that coordinate is zero.
Because the three face edges represent distinct global edges, these arcs
are embedded boundary intervals, mutually disjoint within the attaching
face, even though its vertices coincide in the quotient.

Each cap disc therefore attaches to one base disc along one boundary
interval.  Such an attachment leaves a disc and cannot join two
components.  All cap pieces attach, so none creates an additional
component.  Apply the argument successively to the two caps; the second
original face remains available after the first attachment.  This gives
exactly \(x+b\) discs.

Essentiality also survives these attachments.  In the abstract attaching
disc, each old normal arc is replaced on the complementary boundary disc
by the remainder of the same cap disc's boundary.  The replacement arcs
are disjoint and have exactly the same endpoints and pairing.
Two disjoint arc systems in discs with the same boundary pairing are
related by a PL homeomorphism fixed on the boundary: cut successively
along outermost arcs and extend over the resulting discs.  This
homeomorphism respects every quotient identification on the outer
boundary and hence descends to a homeomorphism of the boundary tori,
identity off the replaced face.  It carries old boundary curves to new
ones.  Thus meridian boundaries remain essential and link boundaries
remain inessential.  This argument uses abstract discs and their
identified boundaries, rather than assuming embedded face closures.

The triangles counted by \(u,v\) are disjoint links of the two new
boundary vertices.  Adding them contributes \(u+v\) inessential discs.
Finally put \(b=h\), \(u=v=0\) for the canonical statement.
\end{proof}

\begin{corollary}[Essential standard rays beyond Q-rays]
\label{fib:non-q-discs}
All seven nonlink rays are discs.  The four with \(x=1\) are essential;
the three with \(x=0\), including \((0,1,1)\), are inessential.
Only \((1,0,0)\), \((0,1,0)\), and \((0,0,1)\) are quadrilateral
extreme directions.  Consequently three of the four essential standard
rays are absent from the Q-extreme-ray list.
\end{corollary}

\subsection{Exact work performed by the full equality arrangement}
\label{fib:arrangement}

The counts below concern the standard phase of \texttt{sector\_rays},
using the \texttt{arrangement} method without a work cap.  It considers
every pair of distinct projected hyperplanes, deduplicates nonnegative
primitive directions, lifts each retained direction, and tests standard
extremality.  These counts describe work after kernel construction.

\begin{proposition}[Attempted, positive and rejected directions]
\label{fib:legacy-counts}
For this family the full equality arrangement has
\begin{equation}
 H=2n+6,\qquad
 \text{pair attempts}=\binom{2n+6}{2},\qquad
 \text{distinct nonnegative directions}=(n+1)^2+3.
 \label{fib:legacy-numbers}
\end{equation}
Exactly \((n-1)(n+3)\) of these directions are rejected by the
standard-extremality test; seven are emitted.
\end{proposition}
\begin{proof}
The distinct old-vertex potentials in the base-offset gauge are
\[
 \{cx:c\in C_n\}\ \cup\ \{Ax-y,Ax-z\},\qquad
 C_n=\{0,F_2,F_3,\ldots,F_{n+2}\}.
\]
The set \(C_n\) has \(n+2\) elements, because the Fibonacci sequence
is strictly increasing from \(F_2\) onwards.  Pairwise equalities and
quadrilateral coordinate hyperplanes give precisely
\begin{equation}
 x=0,\qquad y=(A-c)x\ (c\in C_n),\qquad
 z=(A-c)x\ (c\in C_n),\qquad y=z.
 \label{fib:all-planes}
\end{equation}
The planes \(y=0,z=0\) already occur at \(c=A\).  All displayed
planes are distinct, proving \(H=2n+6\).  Every pair is independent
in three dimensions, and the explicit loop visits every pair.

For \(x>0\), normalize \(x=1\).  Put
\[
 S_n=\{A-c:c\in C_n,\ c\leq A\}.
\]
Only \(B\) is excluded by \(c\leq A\), so \(|S_n|=n+1\).
The intersections in the nonnegative quadrant are exactly the grid
\((y,z)\in S_n\times S_n\): intersections with \(y=z\) merely
repeat its diagonal entries.  The remaining \(x=0\) directions are
\((0,1,0)\), \((0,0,1)\), and \((0,1,1)\).
All grid directions are integral and primitive, since the last base
quadrilateral is one.  This proves the positive-direction count after
deduplication.

Every grid point satisfies \(0\leq y,z\leq A\), hence \(h=0\).
Among them only the four corners of this square are standard extreme;
every other point is a nontrivial convex combination of corners on the
linear canonical-lift cell.  The three \(x=0\) directions are all
extreme.  The rejected count is therefore
\((n+1)^2-4=(n-1)(n+3)\).  In fact every rejected grid direction is
itself a connected meridian disc, by Theorem~\ref{fib:disc-census};
its rejection concerns extremality rather than topology.
\end{proof}

The minimum construction retains three cells and seven directions for
all \(n\), whereas the equality arrangement performs a growing
quadratic number of distinct positive lifts and extremality checks.
The additional \(n+1\) inactive potentials are nonnegative multiples
of \(x\), dominated by the zero potential.  Source construction,
projection costs, coordinate arithmetic and serialization are separate
from the exact loop counts in~\eqref{fib:legacy-numbers}.

\subsection{Encoding size and finite independent evidence}
\label{fib:encoding}

\begin{proposition}[Input and output bit lengths]
\label{fib:bits}
With an explicit numbered face-pairing table, the family input has
\(\Theta(n\log(n+1))\) bits.  The seven complete primitive normal
vectors have \(\Theta(n^2)\) bits in total.  Their largest coordinate
is \(B=F_{n+2}\).  A specialized implementation of the displayed
formulas can construct and serialize the seven vectors in
\(\Theta(n^2)\) bit operations using ordinary binary addition.
\end{proposition}
\begin{proof}
There are linearly many tetrahedron references, a linear fraction of
which require \(\Theta(\log(n+1))\) bits; all permutations have
constant size.  Each output contains linearly many coordinates, each
with \(O(n)\) bits.  Conversely the base coordinates of a meridian
contain \(F_1,\ldots,F_n\), whose binary lengths sum to
\(\Theta(n^2)\); for example a linear fraction of these entries have
linear bit length.  The cap coordinate \(B\) realizes the maximum.
Generate the Fibonacci list by successive additions and substitute into
the seven formulas.  The additions and the full output require
\(O(n^2)\) bit work, matching the output-size lower bound.
\end{proof}

A succinct family recipe needs only the integer \(n\), but the native
solver receives the expanded triangulation.  The specialized formula
bound is not a bound on all stages of the generic sector implementation.
Likewise exponentially many normal pieces need not be expanded to
classify their components; compressed orbit methods provide the relevant
general setting~\cite{AHT2006}.

The retained \texttt{family\_audit.py} checks \(n=1,2,4,8\).
It verifies 300 prescribed parameter-grid instances against the native
canonical lift, regenerates the seven rays, and independently replays
four complete coverage certificates.  Native compressed component
censuses and their independent checker classify all 28 ray instances
and twelve additional mixed directions.  Regina 7.4 independently
classifies all 28 rays and verifies that each source is a solid torus.
Its complete standard vertex lists contain respectively
\(35,39,85,375\) surfaces; restricting each list to the supplied sector
and removing pure links gives exactly the seven proved rays.
These finite results test the implementation against the formulas and
an external oracle; the all-size conclusions follow from the preceding
proofs.
